\documentclass[11pt,a4paper]{article}
\usepackage[T1]{fontenc}
\usepackage{lmodern,amsmath,amssymb,booktabs,microtype}
\usepackage[margin=24mm]{geometry}
\usepackage[colorlinks=true,urlcolor=blue,linkcolor=blue]{hyperref}
\newcommand{\R}{\mathbb R}
\newcommand{\F}{F_1}
\newcommand{\eps}{\epsilon}
\title{Gaussian pentagon matrix elements:\\a certified heat cutoff and unresolved quadrature}
\author{Alan Day \and Chengheng\thanks{AI research assistant (OpenAI).}}
\date{Working research note --- 1 November 2026}
\begin{document}\maketitle
\begin{abstract}
For the specified Kim--Scarinci phase at $\hbar=1$, a punctured
$H^1$ approximation bounds the heat-cutoff error in a Gaussian
pentagon matrix element by $0.000982$ at $\epsilon=10^{-10}$.
An independent integer/rational enclosure certifies the numerical
constant. Actual FFT and Hermite evaluations give nonzero discrete
matrix elements, but their quadrature errors remain uncertified.
A flat control shows why damping-induced residuals alone cannot
refute the undamped identity. No continuous counterexample or
physical holographic conclusion is established.
\end{abstract}
\section{The precise scope of the certificate}
Use dimensionless coordinates, $d\mu=dx\,dy/(2\pi^2)$ and
\[
 P_j=-\pi i\partial_j,\quad A=\F(X,Y),\quad
 B=\F(P_y,P_x),\quad M=e^{iXY/\pi},
\]
where $\F$ is the bounded paired phase of KS \cite{ks}, defined
almost everywhere by
\[
 \F(x,y)=\Phi^i(x+iy)/\overline{\Phi^i(x+iy)} .
\]
The identity being reproduced is $D=AB-BM^{-1}BMA=0$.
Let $H_\eps=e^{-\eps(P_x^2+P_y^2)}$, $B_\eps=BH_\eps$, and
$D_\eps=AB_\eps-B_\eps M^{-1}B_\eps MA$.
The bounded Fourier-kernel formulation uses the Gaussian-group
framework \cite{gk}; it does not split unbounded chiral factors.

For $f=\sqrt2\,e^{-(x^2+y^2)/(2\pi)}$, normalized in $L^2(\mu)$,
put $v=MAf$ and $E_\eps(u)=\|(1-H_\eps)u\|$.
The preceding bounded-word estimate gives
\[
 |\langle f,(D-D_\eps)f\rangle|
 \leq 2E_\eps(f)+E_\eps(v),\qquad
 E_\eps(f)^2=\frac{2(\pi\eps)^2}{(1+\pi\eps)(1+2\pi\eps)}. \tag{1}
\]
Here we bound the last term explicitly.
All constants are for $\hbar=1$ and this fixed test.
The certificate concerns (1), not numerical integration of
$\langle f,D_\eps f\rangle$.

\newpage
\section{A derivative bound on a punctured box}
Set $q=e^{-\pi}$ and $a_j=(2j+1)\pi$. The known product formula gives
\[
 \F(x,y)=\prod_{j\geq0}g_j(x,y)g_j(y,x),\qquad
 g_j(x,y)=\frac{1+e^{x-a_j-iy}}{1+e^{x-a_j+iy}}. \tag{2}
\]
Each factor has modulus one. For $z=e^{x-a_j+iy}$ a local phase obeys
\[
 |\nabla\arg g_j|=2\left|\frac{z}{1+z}\right|.
\]
Write $u=x-a_j$ and let $\tau\in[-\pi,\pi]$ be the angular displacement
from an odd multiple of $\pi$. Then
\begin{align*}
 \left|\frac{z}{1+z}\right|^2
 &=\frac{e^u}{2(\cosh u-\cos\tau)}
 \leq 1+\frac1{\cosh u-\cos\tau}\\
 &\leq1+\frac{\pi^2}{2(u^2+\tau^2)}. \tag{3}
\end{align*}
The last step uses $\cosh u-1\geq u^2/2$ and
$1-\cos\tau\geq2\tau^2/\pi^2$.

On $[-9,9]^2$, all $j\geq1$ terms are uniformly controlled by
\[
 T_{\rm prod}=
 \frac{e^{9-3\pi}}{(1-e^{9-3\pi})(1-e^{-2\pi})}.
\]
Their radial moduli are less than one, so the geometric denominator
estimate applies term by term. Let $d_x$ be the distance to the
nearest $(\pi,(2k+1)\pi)$, and $d_y$ to the nearest
$((2k+1)\pi,\pi)$. A local phase $\theta=\arg(MF_1)$ satisfies
\[
 |\nabla\theta|\leq C+\frac{\sqrt2\pi}{d_x}
                         +\frac{\sqrt2\pi}{d_y},\qquad
 C=\frac{9\sqrt2}{\pi}+4+4T_{\rm prod}.
\]
Consequently
\[
 |\nabla\theta|^2\leq
 3C^2+6\pi^2(d_x^{-2}+d_y^{-2}). \tag{4}
\]
Only the odd ordinates $\pm\pi,\pm3\pi$ are needed on this box.
These estimates remove more points than are singular in the combined
phase, which is harmless for an upper bound. In particular, no
uncancelled singularity at $(\pi,\pi)$ is asserted.

\newpage
\section{The punctured $H^1$ approximation}
Let $a(s)$ equal one on $[-8,8]^2$, taper linearly in each coordinate
to zero at $|x|=9$ or $|y|=9$, and vanish outside.
Use holes of radius $\rho=10^{-4}$ around the seven points
\[
 {\cal P}=\pi\{(1,-3),(1,-1),(1,1),(1,3),(-3,1),(-1,1),(3,1)\}.
\]
At each point take a radial cutoff that is zero up to $\rho$,
linear up to one at $2\rho$, and one thereafter.
Their transition annuli are disjoint. Let $b$ be their product
and $w=abv$. It is compactly supported and in $H^1$.
The probability density of $|v|^2d\mu$ in ordinary coordinates is
$\varrho(s)=\pi^{-2}e^{-|s|^2/\pi}$. Define
\[
 K(d,R)=\pi^{-2}e^{-(d-R)^2/\pi},\quad
 T=\frac14 e^{-64/\pi},\quad
 J_\rho=\sum_{p\in{\cal P}}4\pi\rho^2K(|p|,2\rho).
\]
The Gaussian union-tail bound gives mass outside $[-8,8]^2$
at most $T$, so $\|v-w\|\leq\sqrt{T+J_\rho}$.
Indeed, bound $\int_L^\infty e^{-x^2/\pi}dx$ by
$\int_L^\infty(x/L)e^{-x^2/\pi}dx$; see also \cite{dlmf}.

For $R=1/2$, $c=3/5$, $d=|p|$, the inverse-distance integral
outside the hole is bounded by
\[
 I_p=
 2\pi K(d,R)\log(R/\rho)
 +\frac1{(1-c)^2d^2}+\frac{e^{-c^2d^2/\pi}}{R^2}. \tag{5}
\]
The first term integrates $\rho<|s-p|<R$.
For the remainder split at $|s|=cd$: inside, $|s-p|\geq(1-c)d$;
outside, the radial Gaussian tail is $e^{-c^2d^2/\pi}$.
Let $\mathcal P'$ be the list $\mathcal P$ with $(\pi,\pi)$ repeated
once, accounting for both product orientations.
Equation (4) and the real Gaussian amplitude derivative give
\[
 G=
 \sqrt{\frac2\pi+3C^2+6\pi^2\sum_{p\in{\cal P}'}I_p}
 +\sqrt{2T}
 +\sqrt{\sum_{p\in{\cal P}}4\pi K(|p|,2\rho)}
 \ \geq\ \|\nabla w\|. \tag{6}
\]
The last two terms bound the square cutoff and hole cutoffs.
Since $\|1-H_\eps\|\leq1$ and $(1-e^{-z})^2\leq z$ for $z\geq0$,
\[
 E_\eps(v)\leq\sqrt{T+J_\rho}+\pi\sqrt\eps\,G. \tag{7}
\]
Thus the original phase need not lie in $H^1$.
\newpage
\section{Elementary enclosure and a direct left-word integral}
Combining (1) and (7) gives
\[
 {\cal B}_\eps=
 2\sqrt{\frac{2(\pi\eps)^2}{(1+\pi\eps)(1+2\pi\eps)}}
 +\sqrt{T+J_\rho}+\pi\sqrt\eps\,G. \tag{8}
\]
Two independent evaluations are saved: 40/70-digit
\texttt{mpmath.iv}, and a Python integer/Fraction implementation.
Since interval support in mpmath is documented as experimental
\cite{mpmath}, the final strict inequalities also use the latter.

The rational implementation rounds outward to 120-bit dyadic
endpoints after arithmetic. It uses the alternating Machin formula
$\pi=16\arctan(1/5)-4\arctan(1/239)$ with next-term tails.
For $\exp(x)$, range reduction gives $0\leq x\leq1/8$; a degree-24
Taylor sum has tail at most its next term divided by $1-x/26$.
Negative arguments use reciprocals.
For $\log y$, powers of two reduce $y$ to $[1,2]$, then
\[
 \log y=2\sum_{j=0}^{39}\frac{z^{2j+1}}{2j+1}+R,\quad
 z=\frac{y-1}{y+1},\quad
 0\leq R\leq\frac{2z^{81}}{81(1-z^2)}.
\]
Integer square roots enclose square roots. The final comparisons
are exact rational endpoint comparisons:
\begin{center}
\begin{tabular}{rrr}
\toprule
$\eps$ & approximate $\mathcal B_\eps$ & certified strict upper\\
\midrule
$10^{-8}$ & .009631522192 & .01\\
$10^{-10}$ & .000981734116 & .000982\\
$10^{-12}$ & .000116763225 & .000117\\
\bottomrule
\end{tabular}
\end{center}
This is an arithmetic certificate for the stated bound, not a
proof-assistant verification of all software or of the pentagon.

An independent continuous integral for the left word is
\[
 \langle f,AB_\eps f\rangle=
 \frac{2}{1+2\pi\eps}\,
 \mathbb E\!\left[
 F_1(S)F_1(T)
 e^{i(S_xT_y+S_yT_x)/\pi}\right], \tag{9}
\]
where $S,T$ are independent centered planar Gaussians with
covariances $\pi I$ and $\pi I/(1+2\pi\eps)$.
This follows directly from the Fourier transform of $f$.
Four-dimensional tensor Gaussian quadrature is evaluated through
two separable matrix products, at orders 128--768.
It uses neither periodic FFT nor a truncated Hilbert basis.
For $F_1=1$, (9) is exactly $(1+\pi\eps)^{-1}$.

\newpage
\section{Actual regulated matrix elements and controls}
We evaluated the same $D_\eps$ with periodic FFT and oversampled
Hermite projection. The heat multiplier is included at both $B$
occurrences. Heat-induced norm loss is not counted as unitary
projection leakage. Thirty-seven records retain both complex word
values, their difference, norms, parameters and scalar product tails.

For $\hbar=1$ and $\eps=10^{-10}$, representative defects are:
\begin{center}
\begin{tabular}{lrr}
\toprule
method/settings & real part & imaginary part\\
\midrule
FFT $1024,\ L=20$ & .00391212915 & -.00431820856\\
FFT $2048,\ L=20$ & .00391230420 & -.00431827530\\
FFT $2048,\ L=28$ & .00396855893 & -.00434478257\\
Hermite $64,\ q=256$ & .00391447296 & -.00424922552\\
Hermite $128,\ q=512$ & .00390554046 & -.00415910020\\
Hermite $256,\ q=768$ & .00392158989 & -.00415995798\\
\bottomrule
\end{tabular}
\end{center}
Their moduli are around $.0057$--$.0059$. At fixed $L$, doubling
FFT size changes the answer much less than changing $L$.
It primarily extends the momentum range while its momentum spacing
stays $\pi^2/L$; apparent convergence in size alone cannot control
periodization and momentum quadrature.

For the left word alone, direct order-768 quadrature at $10^{-10}$
gives
\[
 .914113080399+ .000104730439\,i.
\]
FFT $2048,L=20$ gives $.914141192384+.000103846566\,i$, and
Hermite $256,q=768$ gives $.914061270561+.000098283811\,i$.
These differences have not been enclosed by rigorous error bounds.
Direct quadrature with $F=1$ agrees with the exact Gaussian formula
within $10^{-12}$, and its flat left-word control agrees with FFT
within $10^{-12}$; these controls do not certify the singular
paired-phase integrals.

\paragraph{The flat control is essential.}
For $F_0(x,y)=\exp[-iy\log(1+e^x)/\pi]$, the undamped identity is
known \cite{ks}. FFT and Hermite give undamped scalar residuals
below $10^{-12}$. Yet at $\eps=.01$ both give a defect
approximately $.0334367989605$.
At $\eps=10^{-10}$ it decreases to $3.6133\,10^{-10}$.
Damping does not preserve the exact pentagon. A nonzero regulated
defect without a removal-error estimate proves nothing about $D$.

\newpage
\section{What is established and what is missing}
The heat-removal part of the Gaussian matrix-element error is now
explicitly bounded. At $\eps=10^{-10}$ it is below $.000982$,
including puncture, spatial-tail, derivative and arithmetic
contributions. This improves the earlier empirical phase-loss
estimate without assuming $H^1$ regularity for $MAf$.

The actual regulated words have also been evaluated in two
discretizations, and their left word has an independent
continuous-kernel quadrature. The results motivate a more precise
test, but none supplies a certified enclosure for
$\langle f,D_\eps f\rangle$. It is invalid to subtract $.000982$
from the modulus of a tabulated approximation and call the result
a continuous lower bound.

The next required calculation is a rigorous enclosure
\[
 \langle f,D_\eps f\rangle\in {\cal I}_\eps .
\]
It must account for scalar product truncation, floating-point
rounding, Fourier/physical tails, phase neighborhoods and quadrature.
A true nonzero witness would require
\[
 \inf_{z\in{\cal I}_\eps}|z|>{\cal B}_\eps .
\]
Failure to meet this inequality is inconclusive, not proof of the
identity. The current interval is missing.

Prioritize a singularity-aware enclosure of the right word, which
remains the more difficult integral. The flat and unit controls,
actual complex word values and failed precision targets must remain
visible. Further unbounded-factor analytic continuation is not
needed for this bounded test. Literature may still provide an exact
kernel that simplifies the integration.

The stage contains 37 matrix rows, 12 direct-left rows, six interval
evaluations and three rational certificates. The 188 replay checks
verify saved arithmetic, controls, hashes and exact certificate
regeneration. They do not certify the discrete matrix values.
This is a model calculation in a fixed Hilbert representation,
not an observation, a complete finite-$G$ construction or a no-go
result for de Sitter holography.

\begin{thebibliography}{9}\small
\bibitem{ks} H.~K.~Kim and C.~Scarinci,
\emph{A quantization of moduli spaces of 3-dimensional gravity},
\href{https://arxiv.org/abs/2112.13329v2}{2112.13329v2},
\S6.3 and Proposition~6.15.
\bibitem{gk} S.~Garoufalidis and R.~Kashaev,
\emph{A meromorphic extension of the 3D index},
\href{https://arxiv.org/abs/1706.08132v2}{1706.08132v2},
Appendix~B. The preceding paired-kernel note gives the full dictionary.
\bibitem{dlmf} NIST Digital Library of Mathematical Functions,
\href{https://dlmf.nist.gov/7.8}{\S7.8, Gaussian tail inequalities}.
We use the elementary weaker bound derived in the text.
\bibitem{mpmath} mpmath 1.3.0 documentation,
\href{https://mpmath.readthedocs.io/en/1.3.0/contexts.html}
{Contexts: interval arithmetic}, accessed 1 November 2026.
\end{thebibliography}
\end{document}
